\section{Quantitative compactness without a compactness axiom}
\label{sec:quantitative-analysis}

This section supplies a reusable replacement for three recurring analytical
steps: taking an extremal value on a compact set, obtaining an actual point
from approximate solutions, and inferring local continuity from convexity.
The input data and the precision schedules are part of the theorems.  The
results do not infer a primitive-recursive modulus from arbitrary qualitative
continuity.  The distinction matters in the calculus and Taylor--Cauchy
packages, which explicitly supply their continuity and derivative contracts
\cite{repo-calculus,repo-taylor}.

\subsection{The presentation convention}

For definiteness work on
\[
 K_d=[0,1]^d\cap\Q^d,\qquad
 d_\infty(x,y)=\max_{i<d}|x_i-y_i|.
\]
For $d=0$ the domain contains one empty tuple and its distance is zero.
The same arguments apply to other rational domains when finite nets and
their rounding maps are supplied.  A scalar presentation consists of fixed
primitive-recursive rational functions
\[
 A(x,p),\qquad \mu(k),
\]
with the following two contracts, for rational points in $K_d$:
\begin{align}
 |A(x,p)-A(x,q)|&\le 2^{-p}+2^{-q},
       \label{eq:analysis-coherence}\\
 d_\infty(x,y)\le2^{-\mu(k)}
 &\ \Longrightarrow\ 
 |A(x,p)-A(y,q)|\le2^{-k}+2^{-p}+2^{-q}.
       \label{eq:analysis-continuity}
\end{align}
We write $f(x)$ for this coherent rational name.  Expressions involving
$f(x)$ are notation for rational inequalities at arbitrary requested
precisions, rather than a new sort of completed real numbers in \iPRA.
For example, $f(x)\le r$ means
$A(x,p)\le r+2^{-p}$ for every $p$.  For sums and products one first
allocates the finite error budget and then obtains an equivalent family of
rational inequalities.

Whenever we remove an arbitrarily small positive error from a rational
inequality, the operation has an elementary finite explanation.  If a
positive rational gap is $u/v$, with positive integers $u,v$, a dyadic error
smaller than $1/v$ contradicts that gap.  A suitable exponent is obtained by
bounded search below $v+1$.  Thus arguments such as ``for every precision,
therefore the limiting inequality holds'' can be translated into a
particular rational contradiction at a computable precision.  They are not
applications of a completeness principle.

The functions $A$ and $\mu$ are fixed syntax in the object theory, or are
specific functions already present in a definitional extension.  The
theorem is a schema over such presentations.  In particular, there is no
primitive-recursive universal evaluator for arbitrary primitive-recursive
program codes hidden in the phrase ``evaluate $A$''.

\subsection{A finite-net extremum theorem}

\begin{theorem}[A primitive-recursive infimum name]
\label{thm:finite-net-infimum}
From \eqref{eq:analysis-coherence}--\eqref{eq:analysis-continuity} one can
construct a primitive-recursive sequence of rational points $a_k\in K_d$
and nested rational intervals $[L^*_k,U^*_k]$ satisfying
\begin{align}
 0\le U^*_k-L^*_k&\le\tfrac34\,2^{-k},
       \label{eq:infimum-width}\\
 L^*_k&\le f(x)\quad(x\in K_d),
       \label{eq:infimum-lower}\\
 f(a_k)&\le V+\tfrac34\,2^{-k},
       \label{eq:approx-minimizer}
\end{align}
where $V$ is the name given by these nested intervals.  The name $V$ is a
lower bound for all $f(x)$ and is the unique named value with this lower
bound and approximate-attainment property.

The construction uses
\[
 n_k=\mu(k+2),\qquad
 N_k=(2^{n_k}+1)^d
\]
grid evaluations, all at precision $k+2$.  It does not assert that the
points $a_k$ form a Cauchy name.
\end{theorem}

\begin{proof}
Let $G_n$ be the full dyadic grid whose coordinates are
$0,2^{-n},\ldots,1$.  Enumerate it in a fixed lexicographic order.
For $x\in K_d$, rounding each coordinate down gives a grid point
$\pi_n(x)$ with
\[
 d_\infty(x,\pi_n(x))\le2^{-n}.
\]
Both the enumeration and rounding are rational, bounded computations.

Put $\varepsilon_k=2^{-k}$.  Evaluate every $g\in G_{n_k}$ at
precision $k+2$, let $m_k$ be the least sampled value, and let $a_k$ be
the first point attaining it.  Rational order makes the tie-breaking
decidable.  For every rational $x\in K_d$,
\begin{align*}
 m_k
 &\le A(\pi_{n_k}(x),k+2)\\
 &\le f(\pi_{n_k}(x))+\tfrac14\varepsilon_k\\
 &\le f(x)+\tfrac12\varepsilon_k.
\end{align*}
Also
$f(a_k)\le m_k+\varepsilon_k/4$.
Define
\[
 L_k=m_k-\tfrac12\varepsilon_k,
 \qquad U_k=m_k+\tfrac14\varepsilon_k.
\]
Consequently $L_k\le f(x)$ for every rational $x$, whereas
$f(a_k)\le U_k$.  In particular, for every $i,j$,
\[
 L_i\le f(a_j)\le U_j.
\]
The endpoints are rational, so removing the arbitrarily small evaluation
errors gives the literal rational inequality $L_i\le U_j$.

Now take finite running intersections:
\[
 L^*_k=\max_{i\le k}L_i,
 \qquad U^*_k=\min_{i\le k}U_i.
\]
The cross inequalities just proved show $L^*_k\le U^*_k$.  The
intervals are nested, and their width is at most
$U_k-L_k=3\varepsilon_k/4$.  This proves
\eqref{eq:infimum-width} and defines the name $V$ without choosing a
limit. Explicitly, a scalar sampler for this name is
$V(k)=(L^*_k+U^*_k)/2$. For $l\ge k$ both midpoints are in the
$k$th interval, so $|V(l)-V(k)|\le3\,2^{-k}/4\le2^{-k}+2^{-l}$.
Thus it is regular in the sense of \eqref{eq:RN}.
All its lower endpoints are lower bounds for $f$. In particular,
$V(q)\le f(x)+3\,2^{-q}/8$; after allowing the sampling error of
$f(x)$, this is the regular-name lower-bound inequality at index $q$.
Moreover,
as $V\ge L_k$ and $f(a_k)\le U_k$,
\[
 f(a_k)\le V+(U_k-L_k)
          =V+\tfrac34\varepsilon_k.
\]
These relations can each be unfolded at any finite requested precision.

Finally suppose $W$ is another lower-bound name with approximate
attainment.  Since $W\le f(a_k)\le V+3\varepsilon_k/4$ for every
$k$, we have $W\le V$ as a named inequality.  Apply the corresponding
approximate-attainment property of $W$ to obtain $V\le W$.  No
extremizing point or subsequence is used in this comparison.
\end{proof}

The maximum is obtained by applying the construction to $-f$.  A
rational box other than $[0,1]^d$ is handled by an affine change of
variables with its explicit distance factor.  More generally, only a
primitive-recursive finite net, a rounding map, and a modulus were used.
Thus the result is a general finite-net theorem rather than a special
property of cubes.

The construction resolves a useful part of an appeal to the extreme-value
theorem: a supremum or infimum \emph{value}, together with points as close
to the value as requested, has a direct finite realization.  It does not
resolve the separate task of returning a coherent extremizing point.
For example, choosing alternating points $0,1$ for the constant function
returns exact minimizers at every stage, but those choices are not a
Cauchy sequence.  A rule that identifies approximate solutions at nearby
precisions is necessary for the next step.

\subsection{Quantitative uniqueness makes approximate witnesses coherent}

\begin{theorem}[Coherent extraction from approximate minima]
\label{thm:coherent-minimizer}
In the setting of Theorem~\ref{thm:finite-net-infimum}, suppose a fixed
primitive-recursive function $u$ satisfies the quantitative uniqueness
contract
\begin{equation}
 \left.
 \begin{array}{c}
 f(x)\le V+2^{-u(k)}\\
 f(y)\le V+2^{-u(k)}
 \end{array}\right\}
 \quad\Longrightarrow\quad d_\infty(x,y)\le2^{-k}.
 \label{eq:quantitative-uniqueness}
\end{equation}
Set $v(k)=k+\max_{j\le k}u(j)$ and $x_k=a_{v(k)}$.
Then
\[
 d_\infty(x_k,x_l)\le2^{-\min(k,l)}
                    \le2^{-k}+2^{-l}.
\]
The $x_k$ therefore give a primitive-recursive point name $X$ in the
closed cube.  Its canonical evaluation satisfies $f(X)=V$, and it is
the unique point name with that property.
\end{theorem}

\begin{proof}
If $k\le l$, both $v(k)$ and $v(l)$ are at least $u(k)$.
Equation~\eqref{eq:approx-minimizer} shows that $x_k,x_l$ both meet
the hypotheses of \eqref{eq:quantitative-uniqueness} at precision $k$.
This proves the displayed Cauchy estimate.  All coordinates of all
approximants lie between zero and one, so the same closed inequalities
hold of the point name.

For completeness, a canonical evaluation procedure is explicit.  Replace
$\mu$ by its monotone majorant if necessary and put
\[
 r(p)=p+2+\max_{j\le p+2}\mu(j),\qquad
 B(p)=A(x_{r(p)},p+2).
\]
The distance bound between $x_{r(p)}$ and $x_{r(q)}$, followed by
\eqref{eq:analysis-continuity}, proves coherence of $B$ with ample
dyadic slack.  This defines $f(X)$.  Further,
$v(r(p))\ge r(p)\ge p+2$, so the optimum gap of $x_{r(p)}$ is at
most $2^{-p-2}$.  The evaluation error is also $2^{-p-2}$.
Hence $B(p)$ differs from $V$ by at most $2^{-p-1}$, interpreted using
the interval name for $V$.  This proves equality.

If another point name $Y$ has value $V$, choose rational approximants
to $X,Y$ fine enough for the continuity modulus and for the tolerance in
\eqref{eq:quantitative-uniqueness}.  The equality of the values supplies
the required approximate-minimum premises.  The uniqueness contract and
the coordinate errors show $d_\infty(X,Y)$ is below any prescribed
dyadic tolerance.  Thus $X=Y$ as point names.
\end{proof}

\begin{corollary}[A complete strongly convex minimization scheme]
\label{cor:strong-convexity}
Suppose a positive rational $\lambda$ is supplied and
\begin{equation}
 f\!\left(\frac{x+y}{2}\right)
 \le\frac{f(x)+f(y)}2
       -\frac{\lambda}{8}\,d_\infty(x,y)^2
 \quad(x,y\in K_d).
 \label{eq:midpoint-drop}
\end{equation}
Let $c$ be a natural number with $\lambda2^c\ge8$.
Then $u(k)=2k+c$ satisfies \eqref{eq:quantitative-uniqueness}.
In particular the simpler sequence
\[
 x_k=a_{2k+c}
\]
already gives a coherent name of the unique minimizer.
\end{corollary}

\begin{proof}
The midpoint belongs to $K_d$, so the lower-bound property of $V$
and \eqref{eq:midpoint-drop} give
\[
 V\le f\!\left(\frac{x+y}{2}\right)
   \le V+2^{-u(k)}
                 -\frac{\lambda}{8}d_\infty(x,y)^2.
\]
Thus
\[
 d_\infty(x,y)^2
 \le\frac8\lambda\,2^{-2k-c}\le2^{-2k}.
\]
The distance is rational and nonnegative.  Rational order therefore
gives $d_\infty(x,y)\le2^{-k}$; a positive contrary gap would remain
positive after squaring.  This is precisely the uniqueness contract.
The schedule $2k+c$ is already monotone, and the approximate minima
on this schedule have a vanishing value gap, so the preceding proof
applies directly.

The integer $c$ need not be found by unbounded logarithmic search.
Write $\lambda=a/b$ with positive integers $a,b$.  Searching
$0\le c\le8b$ suffices, since $2^{8b}\ge8b\ge8b/a$.
Consequently every step of the stated schedule is primitive recursive.
\end{proof}

The usual Euclidean $\lambda$-strong-convexity inequality implies
\eqref{eq:midpoint-drop}, because
$\|x-y\|_2^2\ge d_\infty(x,y)^2$.  One may also prove
\eqref{eq:midpoint-drop} directly for a particular rationally presented
objective, without introducing derivatives or Hessians.  For example,
$f(x)=\sum_{i<d}(x_i-b_i)^2$ has $\lambda=2$, by expansion of squares.
The theorem then supplies a complete construction of the minimizing
point, including boundary cases, from finite rational minimization.
This is stronger than merely replacing an extremum by a finite
certificate: it constructs the coherent point and proves its optimality.

For equations one can use the same pattern with a nonnegative residual
$\rho$.  A finite algorithm producing $\rho(a_k)\le2^{-k}$ supplies
approximate witnesses.  A quantitative isolation contract
\[
 \rho(x),\rho(y)\le2^{-u(k)}
 \ \Longrightarrow\ d(x,y)\le2^{-k}
\]
supplies coherence.  Neither ingredient follows from the other.  In
particular, multiple roots may violate this isolation contract.  The
existing FTA package instead controls the movement of its descent
sequence and already distinguishes that argument from unrelated
approximate zeros \cite{repo-fta}.  Our theorem does not replace its
multiple-root construction by an unjustified uniqueness assumption.

\subsection{Removing the extra candidate modulus in Bohr--Mollerup}
\label{sec:bm-modulus-repair}

The Bohr--Mollerup report asks whether the primitive-recursive local
modulus of the candidate can be derived rather than supplied
\cite{repo-bm}.  For its rationally presented functions, the answer is
yes.  In fact, recurrence is unnecessary for this part.

\begin{lemma}[Positive multiplicative convexity implies ordinary convexity]
\label{lem:positive-convexity}
Let $F$ assign strongly positive regular names to positive rationals.
Here strong positivity includes a fixed primitive recursive witness
function giving a positive rational lower bound at each rational
argument, as in the repository's presentation; it is not merely an
unwitnessed assertion of pointwise positivity.
Suppose that, for $s\ge1$, $0\le j\le s$, and
$z=((s-j)x+jy)/s$,
\begin{equation}
 F(z)^s\le F(x)^{s-j}F(y)^j.
 \label{eq:rlc-framework}
\end{equation}
Then
\[
 F(z)\le\frac{s-j}{s}F(x)+\frac{j}{s}F(y).
\]
The implication has a finite arithmetic proof with
primitive-recursive precision demands.
\end{lemma}

\begin{proof}
Finite weighted arithmetic--geometric mean gives
\[
 F(x)^{s-j}F(y)^j
 \le\left(\frac{(s-j)F(x)+jF(y)}s\right)^s.
\]
This is the finite AM--GM theorem already developed in the same
Bohr--Mollerup package, applied to a list with $s-j$ copies of $F(x)$
and $j$ copies of $F(y)$.  Its elementary argument can be recalled
briefly.  For two positive entries, $(u-v)^2\ge0$ proves the result.
For a list of power-of-two length, repeatedly split it into two equal
blocks.  Applying the same inequality to the means of the two blocks,
and raising it to their common length, proves the next level of the
binary product tree.  For a list of arbitrary length $s$, let $m$ be
the least power of two at least $s$, and let $a$ be the mean of the
original list.  Pad the list by $m-s$ copies of $a$.  Its mean is still
$a$, so the power-of-two case gives
$\bigl(\prod_{i<s}u_i\bigr)a^{m-s}\le a^m$.
Cancellation of the positive factor $a^{m-s}$ proves the required
$s$-entry inequality.  There are at most $2s$ leaves in the binary
product tree.  This padding construction avoids a separate downward
induction through all intervening lengths.  At every node the
invariant is an inequality of finite products and rational means;
its finite arithmetic version is proved by a Boolean prefix invariant
over the tree in postorder.

Combine AM--GM with \eqref{eq:rlc-framework} and cancel the common
positive power.  This last step does not decide the order of two real
names.  To verify the desired order at tolerance $\varepsilon>0$,
a putative contrary rational gap of that size forces a positive gap
of at least $\varepsilon^s$ between the $s$th powers, by the finite
factorization of $u^s-v^s$ or the binomial theorem.  Requesting the
power evaluations to error less than half that gap gives a rational
contradiction.  The magnitude bounds needed to evaluate the powers
come from coarse input intervals; $s$, $\varepsilon$, and those
bounds give a primitive-recursive precision demand.  Thus this is
a positive-power monotonicity argument, not real trichotomy.

Here is one explicit demand that also explains the passage from
rational AM--GM to names.  Suppose all relevant values lie in
$[\ell,B]$ with positive rational $\ell$ and integer $B\ge1$.
Approximating each value to error $\eta\le\min(1,\ell/2)$ changes
a product of $s$ factors, and also an $s$th power of their mean,
by at most $s(B+1)^{s-1}\eta$, by telescoping the product.
For a target order tolerance $\varepsilon\le1$, it therefore
suffices to choose a dyadic $\eta$ with
\[
 \eta\le\min\left(1,\frac\ell2,
       \frac{\varepsilon^s}{8s(B+1)^{s-1}}\right).
\]
All quantities on the right are positive rationals.  A dyadic
satisfying the bound is found by a search bounded by their encoded
denominators.  This supplies a concrete primitive-recursive
precision demand for both finite products and the positive-power
comparison; strong positivity supplies $\ell$ for the finitely
many values used in this instance.
\end{proof}

\begin{theorem}[A local Lipschitz bound from two endpoint samples]
\label{thm:convex-local-modulus}
Assume the hypotheses of Lemma~\ref{lem:positive-convexity}.  Fix
positive rationals $\alpha\le\beta$.  Put
\[
 a=\frac\alpha2,\qquad b=\beta+1,\qquad
 \delta=\min\left(\frac\alpha2,1\right).
\]
Let $M\ge1$ be an integer upper bound for $F(a)$ and $F(b)$,
obtained from their precision-zero intervals.  Then, for rational
$x,y\in[\alpha,\beta]$,
\begin{equation}
 |F(x)-F(y)|\le\frac M\delta\,|x-y|.
 \label{eq:convex-lipschitz}
\end{equation}
In particular, with $H=\lceil M/\delta\rceil$,
\begin{equation}
 \mu_F(\alpha,\beta,k)=k+H
 \label{eq:explicit-bm-modulus}
\end{equation}
is a primitive-recursive local uniform-continuity modulus.
\end{theorem}

\begin{proof}
The case $x=y$ is immediate; rational order allows us otherwise to
assume $x<y$.  Express $y$ as the rational convex combination of
$x$ and $b$:
\[
 y=\frac{b-y}{b-x}x+\frac{y-x}{b-x}b.
\]
Lemma~\ref{lem:positive-convexity} gives
\begin{align*}
 F(y)-F(x)
 &\le\frac{y-x}{b-x}\bigl(F(b)-F(x)\bigr)\\
 &\le\frac{M}{\delta}(y-x).
\end{align*}
Here positivity gives $F(x)\ge0$, and $b-x\ge\delta$.
Likewise express $x$ as a rational convex combination of $a$ and
$y$:
\[
 x=\frac{y-x}{y-a}a+\frac{x-a}{y-a}y.
\]
The resulting inequality is
\[
 F(x)-F(y)
 \le\frac{y-x}{y-a}\bigl(F(a)-F(y)\bigr)
 \le\frac M\delta(y-x).
\]
Together these are \eqref{eq:convex-lipschitz}.  The denominators
are positive rationals with the displayed lower bound, so all
division is quantitatively justified.

For a Cauchy sampler $A_F$, rather than nested input intervals, one
may take
\[
 M=1+\left\lceil\max\bigl(|A_F(a,0)|,|A_F(b,0)|\bigr)\right\rceil.
\]
Coherence gives $F(a)\le A_F(a,0)+1\le M$ and the same for $b$.
Thus $M$ uses only two fixed-precision evaluations.
Since $H\ge1$ and $H\le2^H$, the hypothesis
$|x-y|\le2^{-(k+H)}$ gives
\[
 |F(x)-F(y)|\le H\,2^{-(k+H)}\le2^{-k}.
\]
This proves \eqref{eq:explicit-bm-modulus}.  All operations are
integer or rational arithmetic, a ceiling, and the two evaluations.
\end{proof}

This removes $M_F$ from the \emph{independent input data} for the
Bohr--Mollerup candidate on positive rationals.  Its strong positivity,
normalization, recurrence, and rational multiplicative convexity now
suffice to derive both its rational uniqueness theorem and the modulus
needed for its canonical extension to regular positive real names.
One can then use the repository's completion argument with the
derived modulus.

The scope of this repair is precise.  A rational presentation determines
its canonical continuous extension.  Assumptions imposed only at rational
arguments do not, by themselves, constrain arbitrary separately specified
values on irrational arguments.  If $F$ is presented on real names from
the outset, one must require that this presentation is the canonical
extension, or prove its compatibility with the rational presentation.
The theorem supplies the previously missing modulus in the stated
rational-presentation setting; it is not an extension theorem for an
unrestricted set-theoretic function.

\subsection{Why metastability does not supply a Cauchy rate}

There is a useful obstruction which can itself be stated with finite
arithmetic.  Let $h_e(n)$ be zero if machine $e$ has not halted in its
first $n$ steps, and one otherwise.  Bounded simulation is primitive
recursive, and $h_e$ is nondecreasing with at most one jump.

Fix a primitive-recursive control function $g$ with $g(n)\ge n$.
For any $e$, one of the two disjoint integer windows
\[
 [0,g(0)]\quad\hbox{and}\quad
 [g(0)+1,g(g(0)+1)]
\]
has no jump.  Indeed, each nonconstant window contains a jump strictly
between its first and last sampled index, and these windows cannot
both contain the unique jump.  Inspecting both finite windows therefore
finds an $N\in\{0,g(0)+1\}$ such that
\[
 |h_e(i)-h_e(j)|=0
 \quad\hbox{for all }i,j\in[N,g(N)].
\]
This is a uniform primitive-recursive metastability procedure for
every requested positive tolerance.

On the other hand, a uniform computable Cauchy rate $R(e)$ for tolerance
$1/2$ would decide halting: simulate $e$ for $R(e)$ steps.  If it has
not halted by then, it cannot halt later, because the values at
$R(e)$ and at the later halting time would differ by one.  Hence no
such computable rate, and in particular no primitive-recursive rate,
exists uniformly in $e$.  This is a metamathematical obstruction; it
does not add a machine-halting principle to \iPRA.

Consequently a finite fluctuation bound, or a theorem giving a good
finite window for every control function, cannot be silently substituted
for the modulus needed to construct a regular name.  In the positive
results above the modulus comes from an explicit continuity or
uniqueness inequality, and its dependence on the requested precision
is displayed.

\subsection{Arithmetic induction and the Lean implementation path}

For Theorem~\ref{thm:finite-net-infimum}, encode rational tuples and
the grid as finite natural-number codes.  The minimum loop maintains
the current smallest sampled rational, its first index, and the
Boolean conjunction of comparisons against all samples visited so
far.  The invariant is quantifier-free after replacing a bounded
conjunction by its primitive-recursive Boolean fold.  The running
maximum and minimum used for interval intersection have the same
kind of invariant.  Grid correctness is proved coordinate by coordinate
using the quotient/remainder identity for rational rounding.  The
continuity contract is instantiated at the one rounded point; one
never inducts on a formula asserting a property of all real points.

For Theorem~\ref{thm:coherent-minimizer}, the finite maximum defining
$v$ is another bounded fold.  Having fixed $k,l$, the proof calls the
uniqueness contract only once, at $\min(k,l)$.  The strong-convexity
corollary uses a midpoint, two value comparisons, and an inequality
of rational squares.  None of these steps uses induction on a
formula quantifying over all names or all precisions.

For Theorem~\ref{thm:convex-local-modulus}, rational convexity is a
theorem schema for each displayed finite pair of arguments and
rational weight.  The finite AM--GM arithmetic has loop bounds at
most $2s$ and finite product-error contracts.  Once ordinary
convexity is available, the modulus proof uses exactly two convexity
instances.  The choice of $M$ is a primitive-recursive expression
in two coarse samples; there is no search for an unknown neighborhood
or a positive real gap.

A useful Lean implementation order is the following.
\begin{enumerate}
\item Define rational codes, tuple lookup, $d_\infty$, dyadic grids,
      rounding, the first-index minimum fold, and running intersections.
      Prove their elementary finite specifications.  At the semantic
      level a structure may collect $A$, $\mu$, coherence and continuity,
      but the object-language theorem uses their fixed PR definitions.
\item Prove the \emph{finite} lemma
      \[
      m_k\le A(x,q)+\tfrac12\,2^{-k}+2^{-q}
      \]
      directly for arbitrary rational $x$ and precision $q$.  Prove
      $A(a_k,q)\le U_k+2^{-q}$ and the cross-interval inequality
      $L_i\le U_j$.  These are the mathematical core of the infimum
      constructor, and give its object-language targets without
      importing Lean's real compactness theorems.
\item Prove a generic rational lemma saying that two approximate
      minimizers under \eqref{eq:midpoint-drop} have squared distance
      at most $8\varepsilon/\lambda$.  Instantiate it with
      $\varepsilon=2^{-2k-c}$ to verify the point constructor.
\item Formalize weighted AM--GM, positive-power order reflection with
      an explicit precision demand, the two secant inequalities, and
      $H2^{-H}\le1$.  The candidate modulus then follows by composition.
\item Provide the arithmetic derivations of these fixed schemas and
      feed them to the common finite-contract compiler.  A Lean theorem
      about ordinary $\mathbb R$ is useful as a semantic check, but
      is not the object-language \iPRA proof.
\end{enumerate}

The contracts above identify the exact remaining formalization work.
They are definitions and proved mathematical specifications, rather
than unexplained assumptions named after the desired conclusion.
